\documentclass[11pt,a4paper]{article}
\usepackage[utf8]{inputenc}
\usepackage[T1]{fontenc}
\usepackage{lmodern}
\usepackage{amsmath}
\usepackage[margin=2cm]{geometry}
\usepackage{enumitem}
\pagestyle{plain}
\setlength{\parindent}{0pt}
\setlength{\parskip}{0pt}
\begin{document}
{\Large\bfseries Tarea de materia y átomo}
\vspace{12pt}
\begin{enumerate}[label=\textbf{\arabic*.},leftmargin=*,itemsep=11pt,topsep=0pt]
\item \begin{minipage}[t]{\linewidth}
¿Cuál de estas es una mezcla heterogénea?

\vspace{5pt}
\begin{tabular}{@{}p{0.47\linewidth}p{0.47\linewidth}@{}}
a)~vinagre & b)~agua con arena \\
c)~agua con sal disuelta & d)~aire limpio \\
\end{tabular}
\end{minipage}

\item \begin{minipage}[t]{\linewidth}
¿Cuál de estos es un cambio químico?

\vspace{5pt}
\begin{tabular}{@{}p{0.47\linewidth}p{0.47\linewidth}@{}}
a)~un clavo que se oxida & b)~azúcar que se disuelve en agua \\
c)~alcohol que se evapora & d)~un vaso que se rompe \\
\end{tabular}
\end{minipage}

\item \begin{minipage}[t]{\linewidth}
¿Cómo se llama el paso de gas a líquido?

\vspace{5pt}
\begin{tabular}{@{}p{0.47\linewidth}p{0.47\linewidth}@{}}
a)~condensación & b)~solidificación \\
c)~sublimación & d)~fusión \\
\end{tabular}
\end{minipage}

\item \begin{minipage}[t]{\linewidth}
El ${}^{31}\mathrm{P}$ tiene 15 protones. ¿Cuántos neutrones tiene?

\vspace{5pt}
\begin{tabular}{@{}p{0.47\linewidth}p{0.47\linewidth}@{}}
a)~31 & b)~16 \\
c)~15 & d)~46 \\
\end{tabular}
\end{minipage}

\item \begin{minipage}[t]{\linewidth}
¿Cuántos protones, neutrones y electrones tiene el ${}^{59}_{27}\mathrm{Co}$ neutro?

\vspace{5pt}
\begin{tabular}{@{}p{0.47\linewidth}p{0.47\linewidth}@{}}
a)~27 p, 59 n y 27 e & b)~27 p, 32 n y 27 e \\
c)~32 p, 27 n y 32 e & d)~27 p, 32 n y 0 e \\
\end{tabular}
\end{minipage}

\item \begin{minipage}[t]{\linewidth}
El ion cloruro, $\mathrm{Cl}^{-}$, tiene 17 protones. ¿Cuántos electrones tiene?

\vspace{5pt}
\begin{tabular}{@{}p{0.47\linewidth}p{0.47\linewidth}@{}}
a)~18 & b)~16 \\
c)~35 & d)~17 \\
\end{tabular}
\end{minipage}

\item \begin{minipage}[t]{\linewidth}
Un ion tiene 12 protones y 10 electrones. ¿Cuál es?

\vspace{5pt}
\begin{tabular}{@{}p{0.47\linewidth}p{0.47\linewidth}@{}}
a)~$\mathrm{Mg}^{2-}$ & b)~$\mathrm{Mg}^{2+}$ \\
c)~$\mathrm{Ne}$ & d)~$\mathrm{Na}^{+}$ \\
\end{tabular}
\end{minipage}

\item \begin{minipage}[t]{\linewidth}
Un átomo neutro tiene 19 protones y 20 neutrones. ¿Cuál es su notación?

\vspace{5pt}
\begin{tabular}{@{}p{0.47\linewidth}p{0.47\linewidth}@{}}
a)~${}^{39}_{19}\mathrm{K}$ & b)~${}^{39}_{20}\mathrm{Ca}$ \\
c)~${}^{20}_{19}\mathrm{K}$ & d)~${}^{19}_{39}\mathrm{K}$ \\
\end{tabular}
\end{minipage}

\item \begin{minipage}[t]{\linewidth}
¿Cuál de estos pares son isótopos?

\vspace{5pt}
\begin{tabular}{@{}p{0.47\linewidth}p{0.47\linewidth}@{}}
a)~${}^{1}\mathrm{H}$ y ${}^{2}\mathrm{H}$ & b)~${}^{23}\mathrm{Na}$ y ${}^{24}\mathrm{Mg}$ \\
c)~${}^{40}\mathrm{Ar}$ y ${}^{40}\mathrm{K}$ & d)~${}^{14}\mathrm{C}$ y ${}^{14}\mathrm{N}$ \\
\end{tabular}
\end{minipage}

\item \begin{minipage}[t]{\linewidth}
¿Quién propuso que los electrones giran en niveles de energía fijos?

\vspace{5pt}
\begin{tabular}{@{}p{0.47\linewidth}p{0.47\linewidth}@{}}
a)~Rutherford & b)~Thomson \\
c)~Dalton & d)~Bohr \\
\end{tabular}
\end{minipage}

\end{enumerate}
\end{document}
